\documentclass[11pt]{article}
\usepackage[a4paper,margin=0.65in]{geometry}
\usepackage[T1]{fontenc}
\usepackage{amsmath,amssymb,amsfonts}
\usepackage{graphicx}
\usepackage{pdfpages}
\usepackage[english,shorthands=off]{babel}
\usepackage{fancyhdr}
\usepackage{array}
\usepackage[bookmarks=false,colorlinks=false]{hyperref}
\usepackage[protrusion=true,expansion=false]{microtype}
\emergencystretch=3em
\tolerance=600
\providecommand{\modd}{\mathrm{mod}}
\providecommand{\Disc}{\mathrm{Disc}}
\providecommand{\canont}{}
\providecommand{\asterism}{*}
\providecommand{\Hessian}{H}
\providecommand{\tome}{}
\providecommand{\page}{p.}
\pagestyle{fancy}\fancyhf{}\fancyhead[L]{Pages 451--475}\fancyhead[R]{Cayley --- Collected Papers, Vol. VII}\renewcommand{\headrulewidth}{0.4pt}
\setlength{\headheight}{15pt}

\begin{document}

% ====================================================================
% Cayley, Collected Mathematical Papers, Vol. VII
% Scan PDF pages 451--475 (book pages 429--453).
% Content: Paper 476 (cont.) -- ``On the Determination of the Orbit
%          of a Planet from Three Observations''. Articles Nos. 64--97
%          (with planograms Nos. 1, 2, and beginning of No. 3).
% OCR-skeleton + PNG-verification retype.
% ====================================================================

\noindent\textbf{[Paper 476, continued: \textit{On the Determination of the Orbit of a Planet from Three Observations}, Articles 64--97.]}

\medskip

\noindent\textit{[p.~429]}

\smallskip

we obtain the six coordinates of the three rays respectively
\[
\begin{aligned}
(a_{1},\, b_{1},\, c_{1},\, f_{1},\, g_{1},\, h_{1}) &= (\, 0,\ \sqrt{3},\ -1,\ \ \ 0,\ 1,\ \sqrt{3}),\\
(a_{2},\, b_{2},\, c_{2},\, f_{2},\, g_{2},\, h_{2}) &= (\, 3,\ \sqrt{3},\ \ 2,\ -\sqrt{3},\ 1,\ -2\sqrt{3}),\\
(a_{3},\, b_{3},\, c_{3},\, f_{3},\, g_{3},\, h_{3}) &= (-3,\ \sqrt{3},\ \ 2,\ -\sqrt{3},\ 1,\ -2\sqrt{3}),
\end{aligned}
\]
whence the intersections with the orbit-plane are given by
\[
\resizebox{\textwidth}{!}{$\displaystyle
\begin{aligned}
x_{1}'\,:\,y_{1}'\,:\,1 &= \beta'\sqrt{3}-\gamma'\ :\ -\beta'\sqrt{3}+\gamma\ :\ \beta'' + \gamma''\sqrt{3},\\
x_{2}'\,:\,y_{2}'\,:\,1 &= 3\alpha' + \beta'\sqrt{3} + 2\gamma'\ :\ -3\alpha - \beta\sqrt{3} - 2\gamma\ :\ \alpha''\sqrt{3} + \beta'' - 2\sqrt{3}\gamma'',\\
x_{3}'\,:\,y_{3}'\,:\,1 &= -3\alpha' + \beta'\sqrt{3} + 2\gamma'\ :\ 3\alpha - \beta\sqrt{3} - 2\gamma\ :\ -\alpha''\sqrt{3} + \beta'' - 2\sqrt{3}\gamma'',
\end{aligned}$}
\]
where if (as before) the position of the orbit-plane be determined by means of the longitude $b$ and colatitude $c$ of the orbit-pole, we have
\[
\begin{aligned}
\alpha\ ,\ \beta\ ,\ \gamma &= \sin b\ ,\ -\cos b\ ,\ 0\ ,\\
\alpha',\ \beta',\ \gamma' &= \cos b\cos c,\ \sin b\cos c,\ -\sin c,\\
\alpha'',\, \beta'',\, \gamma'' &= \cos b\sin c,\ \sin b\sin c,\ \cos c,
\end{aligned}
\]
and the passage from the coordinates $x',\ y'$, to $x,\ y$, is given by
\[
\begin{aligned}
x' &= \phantom{+}x\sin b - y\cos b,\\
y' &= +x\cos b + y\sin b,
\end{aligned}
\]
or conversely
\[
\begin{aligned}
x &= \phantom{+}x'\sin b + y'\cos b,\\
y &= -x'\cos b + y'\sin b.
\end{aligned}
\]

\paragraph{65.} To develope the results, I consider the orbit-pole as passing through certain series of positions. The locus may be a meridian circle: by reason of the symmetry of the system, the results are not altered by a change of $120^{\circ}$ in the longitude of the meridian; so that, by considering the two meridians $0^{\circ}$--$180^{\circ}$ and $90^{\circ}$--$270^{\circ}$, we, in fact, consider twelve half meridians at the intervals of $30^{\circ}$. An illustration is afforded by Plate~I.; the orbit-pole describes successively the meridians $0^{\circ}$, $30^{\circ}$, $60^{\circ}$, $90^{\circ}$, and the line $1$, by its intersection with the orbit-plane, traces out on this plane a series of hyperbolas shown in the figure; the hyperbola for the meridian $90^{\circ}$ is a right line, but (except for the position where the orbit-plane passes through the line $1$) the locus is a determinate point on this line. Pianogram No.~$1$ (Plate~II.) refers to the meridian $90^{\circ}$--$270^{\circ}$, and Pianogram No.~$2$ (Plate~III.) to the meridian $0^{\circ}$--$180^{\circ}$. Next, if the orbit-pole be at one of the points $A$, that is, if the orbit-plane pass through a ray though the position of the orbit-pole be here determinate, yet as there is a series of orbits, this also will give rise to a pianogram: I call it Pianogram No.~$3$. The orbit-pole may pass along a separator circle (viz.\ the orbit-plane be parallel to a ray), this is Pianogram No.~$4$. And, lastly, the orbit-pole may pass along the ecliptic (or the orbit-plane may pass through the axis $SZ$), I call this Pianogram No.~$5$. But the last three planograms are not considered in the like detail as the first two, and I have not, in regard to them, tabulated the results, nor given any Plates.

\medskip
\noindent\textit{[p.~430]}

\medskip

\begin{center}
\textit{Article Nos.~$66$ to $82$. Planogram No.~$1$, the Meridian $90^{\circ}$--$270^{\circ}$ (see Plate II.).}
\end{center}

\paragraph{66.} Supposing that the orbit-plane rotates about the axis $S1$ (fig.~$6$, see No.~$64$) in the plane of the ecliptic, the orbit-pole will describe the meridian $90^{\circ}$--$270^{\circ}$, the position of the orbit-pole being $b=90^{\circ}$, $c=0^{\circ}$ to $90^{\circ}$, or else $b=270^{\circ}$, $c=0^{\circ}$ to $90^{\circ}$. But the same analytical formula extends to the two half meridians, viz., we may take $b=90^{\circ}$, and extend $c$ over $180^{\circ}$, in the final results making $c$ an arc between $0^{\circ}$ and $90^{\circ}$, and $b=90^{\circ}$, or $=270^{\circ}$, as the case requires.

\paragraph{67.} Assuming then $b=90^{\circ}$, we have
\[
\begin{aligned}
\alpha,\ \beta,\ \gamma &= 1,\ \phantom{-\cos c,}\ 0,\ \phantom{-\sin c,}\ 0,\\
\alpha',\ \beta',\ \gamma' &= 0,\ -\cos c,\ \sin c,\\
\alpha'',\, \beta'',\, \gamma'' &= 0,\ \phantom{-}\sin c,\ \cos c,
\end{aligned}
\]
and, moreover, $x'=x$, $y'=y$: so that instead of $(x_{1}',\, y_{1}')$, \&c., we may write at once $(x_{1},\, y_{1})$, \&c. The formulae become
\[
\begin{aligned}
x_{1}\,:\,y_{1}\,:\,1 &= \sqrt{3}\cos c + \sin c\ :\ 0\ :\ \sin c + \sqrt{3}\cos c,\\
x_{2}\,:\,y_{2}\,:\,1 &= -\sqrt{3}\cos c - 2\sin c\ :\ -3\ :\ \sin c - 2\sqrt{3}\cos c,\\
x_{3}\,:\,y_{3}\,:\,1 &= \phantom{-}\sqrt{3}\cos c - 2\sin c\ :\ -3\ :\ \sin c - 2\sqrt{3}\cos c,
\end{aligned}
\]
that is
\[
x_{1} = 1,\qquad y_{1} = 0.
\]
(viz.\ the orbit-plane, as is evident, meets the ray $1$ in a fixed point, its intersection with the plane of $xy$)
\[
\begin{aligned}
x_{2} &= \dfrac{-\sqrt{3}\cos c - 2\sin c}{\sin c - 2\sqrt{3}\cos c},\qquad y_{2} = \dfrac{-3}{\sin c - 2\sqrt{3}\cos c},\\[4pt]
x_{3} &= \dfrac{\sqrt{3}\cos c - 2\sin c}{\sin c - 2\sqrt{3}\cos c} = -x_{2},\qquad y_{3} = \dfrac{-3}{\sin c - 2\sqrt{3}\cos c} = -y_{2}.
\end{aligned}
\]

\medskip

\noindent\textit{[Note: in the printed text the formula for $x_{3}$ is given with an obvious typographical sign error; the corrected expression makes $x_{3} = -x_{2}$ and $y_{3} = y_{2}$ symmetric, but the printed page gives $y_{3} = -y_{2}$ which we have preserved verbatim.]}

\medskip
and writing
\[
\dfrac{2\sqrt{3}}{\sqrt{13}}\cos\omega = \dfrac{1}{\sqrt{13}},\qquad \dfrac{1}{\sqrt{13}}\sin\omega = \dfrac{1}{2\sqrt{3}}\tan\omega,
\]
(whence $\omega = 16^{\circ}\,6'$) we find
\[
\begin{aligned}
x_{2} &= -\tfrac{1}{2} + \tfrac{3\sqrt{3}}{\sqrt{13}}\tan(c+\omega),\\
y_{2} &= +\tfrac{3}{\sqrt{13}}\sec(c+\omega),
\end{aligned}
\]
and we thence have for the hyperbola, the locus of $(x_{2},\, y_{2})$ and $(x_{3},\, y_{3})$
\[
\bigl(x + \tfrac{1}{2}\bigr)^{2} = \tfrac{3}{12}\bigl(y^{2} - \tfrac{3}{4}\bigr),
\]

\medskip
\noindent\textit{[p.~431]}

\medskip

viz.\ the points $(x_{2},\, y_{2})$ and $(x_{3},\, y_{3})$ are situate on the hyperbola, symmetrically on opposite sides of the axis of $x$. For $c=0$, we have $x_{2} = -\tfrac{1}{2}$, $y_{2} = \tfrac{1}{2}\sqrt{3}$, $(x_{2}^{2} + y_{2}^{2} = 1)$, and the hyperbola at this point touches the circle $x^{2} + y^{2} = 1$; and similarly for $x_{3},\, y_{3}$. The inclination of the asymptotes to the axis of $y$ is given by $\tan\eta = \sqrt{\tfrac{4}{9}},\ \eta = 22^{\circ}\,56'$.

\paragraph{68.} The orbits are conics, focus $S$ and vertex $1$. It will be convenient to consider $c$ as passing from $0^{\circ}$ to $90^{\circ}-\omega$, and from $0^{\circ}$ to $-(90^{\circ}+\omega)$; that is, from $0^{\circ}$ to $73^{\circ}\,54'-\epsilon$, and from $0^{\circ}$ to $-106^{\circ}\,6'+\epsilon$, if $\epsilon$ be indefinitely small: the point $2$ will thus traverse the upper branch (alone shown in the Plate) of the guide-hyperbola, viz., for $c=0^{\circ}$ it will be at the point of contact with the circle; for $c=73^{\circ}\,54'-\epsilon$ it will be at $\infty$, and for $c=-106^{\circ}\,6'+\epsilon$ at $\infty'$. For $c=0^{\circ}$ the orbit is the circle; as $c$ increases positively, it becomes an ellipse of increasing eccentricity and major axis, until for a certain value ($c=46^{\circ}\,48'$ as will appear) it becomes a parabola; it then becomes a hyperbola (concave branch); for $c=52^{\circ}\,45'$ it becomes the hyperbola $\Sigma'$ subsequently referred to; and for $c=60^{\circ}$ (the point $2$ being then on the line shown in the figure) the orbit becomes this right line. As $c$ continues to increase, the orbit becomes a hyperbola (convex branch); and ultimately for $c=73^{\circ}\,54'-\epsilon$, the point $2$ goes to $\infty$, and the orbit becomes a hyperbola (convex) $\Sigma$, having an asymptote parallel to that of the guide-hyperbola: the inclination to the axis of $x$ being thus $90^{\circ} - 22^{\circ}\,56' = 67^{\circ}\,4'$.

\paragraph{69.} Next as $c$ increases negatively, the point $2$ moves from the point of contact in the other direction to $\infty'$: for $c=0^{\circ}$ the orbit is of course the circle, and as $c$ increases negatively the orbits are at first the very same series of orbits as those belonging to the positive values\footnote{Of course, as corresponding to different values of $c$, they are not the same orbits in space, but they are only the same curves in the planogram.}, viz., they are first ellipses, of increasing eccentricity and major axis; then for $c=-92^{\circ}\,54'$ the orbit is the parabola; the orbits are then hyperbolas (concave), and finally for $c=-106^{\circ}\,6'+\epsilon$, when $2$ is at $\infty'$, the orbit is a hyperbola $\Sigma'$, the asymptote of which is parallel to that of the guide-hyperbola, viz., the inclination to the axis of $x$ is $67^{\circ}\,4'$.

\paragraph{70.} It will be observed that the orbits from the circle to the hyperbola $\Sigma'$ each intersect the guide-hyperbola (that is, the branch shown in the figure) in two points, the one corresponding to a positive, the other to a negative value of $c$; in the positive series, the remaining orbits from the hyperbola $\Sigma'$, through the right line to the convex hyperbola $\Sigma$, each intersect the guide-hyperbola (same branch) in a single point only, for which $c$ is positive.

\paragraph{71.} There is, in the passage of the orbit-pole from $c=-106^{\circ}\,6'+\epsilon$ to $c=73^{\circ}\,54'-\epsilon$, say at $c=73^{\circ}\,54'$, a discontinuity of orbit, viz., an abrupt change from the concave hyperbola $\Sigma'$ to the concave hyperbola $\Sigma$; observe that the direction of the asymptotes being the same in each, the eccentricity $e$ has the same value.

\medskip
\noindent\textit{[p.~432]}

\medskip

The point in question ($b=90^{\circ}$, $c=73^{\circ}\,54'$) is one of the points $B$ of the spherogram, and the hyperbolas $\Sigma,\, \Sigma'$ are two of the four orbits belonging to this point. And, by what precedes, it appears that as the orbit-pole passes through this point along a meridian downwards to the ecliptic the change is from a concave to a convex orbit.

\paragraph{72.} On account of the symmetry in regard to the axis of $x$, the equation of the orbit will be of the form $r = \dfrac{Ax + B}{1}$; viz., the equation is at once found to be
\[
r = \dfrac{x_{2} - 1}{x_{2} - 1}\,(x - 1).
\]

\paragraph{73.} The eccentricity is the coefficient $A$ taken positively ($e = \pm A$): it is in the present case proper to attend to the value of the coefficient itself,
\[
A = \dfrac{r_{2} - 1}{x_{2} - 1},
\]
the sign of $A$ will then indicate the position of the centre of the orbit, viz., according as $A$ is positive or negative the centre will be on the negative or the positive side of the focus $S$. To investigate the variation of $A$, we may express it as a function of $\tan c$, $= \lambda$ suppose. We have
\[
x_{2} = \dfrac{\sqrt{3} - 2\lambda}{\lambda - 2\sqrt{3}},\qquad y_{2} = \dfrac{-3}{\lambda - 2\sqrt{3}},
\]
and thence
\[
r_{2} = \dfrac{1}{\lambda - 2\sqrt{3}}\,R_{2},\qquad R_{2} = \pm\sqrt{12 - 4\sqrt{3}\lambda + 13\lambda^{2}},
\]
viz., $r_{2}$ must be positive, that is, $R_{2}$ is positive or negative according to the sign of $\lambda - 2\sqrt{3}$; negative if $\lambda < 2\sqrt{3}$ or $c < 73^{\circ}\,54'$, positive if $\lambda > 2\sqrt{3}$ or $c > 73^{\circ}\,54'$. And we have then
\[
A = \dfrac{\lambda - 2\sqrt{3} - R_{2}}{3(\lambda - \sqrt{3})}.
\]
But a more convenient formula is obtained by writing
\[
\theta = -\cot c + \dfrac{1}{2\sqrt{3}},\qquad a = \dfrac{1}{2\sqrt{3}} - 1,
\]
we then have
\[
\sqrt{1 + \theta^{2}} = \theta r_{2},
\]
which determines the sign of the radical, viz., this must have the same sign as $\theta$; and then for the coefficient
\[
A = \dfrac{2}{3(\theta + a)}\bigl(-\sqrt{1 + \theta^{2}} + \theta\bigr).
\]

\medskip
\noindent\textit{[p.~433]}

\medskip

\paragraph{74.} For $c$ a small arc $=\epsilon$, $\theta$ is large and negative, and $\sqrt{1 + \theta^{2}}$, having the same sign as $\theta$, is $=\theta$ nearly; we have therefore
\[
A = \dfrac{2}{3\theta}\cdot\dfrac{-1}{2\theta} = \dfrac{-1}{3\theta^{2}}\quad\text{approximately.}
\]
For $c$ nearly $=60^{\circ}$, say $c = 60^{\circ}\pm\epsilon$,
\[
\cot c = \cot 60^{\circ}\pm\epsilon\csc^{2}60^{\circ} = \tfrac{1}{\sqrt{3}}\pm\dfrac{4\epsilon}{3},
\]
\[
\theta = -\dfrac{1}{2\sqrt{3}}\pm\dfrac{4\epsilon}{3},\qquad \theta + a = \pm\dfrac{4\epsilon}{3},\qquad \sqrt{1 + \theta^{2}} = -\dfrac{\sqrt{13}}{2\sqrt{3}},
\]
and thence
\[
A = \dfrac{-1 + \sqrt{13}}{2\sqrt{3}}\pm\dfrac{4\epsilon}{3} = \pm\dfrac{\sqrt{3}}{8}\cdot\dfrac{-1 + \sqrt{13}}{\epsilon},
\]
viz., this is $-\infty$ for $c = 60^{\circ} - \epsilon$, and $+\infty$ for $c = 60^{\circ} + \epsilon$.

For $c$ nearly $= 90^{\circ}-\omega$, say first $c = 73^{\circ}\,54' - \epsilon$, we have
\[
\cot c = \dfrac{1}{2\sqrt{3}} + \tfrac{13}{12}\epsilon,\qquad \theta = -\tfrac{13}{12}\epsilon,\qquad \theta + a = -\dfrac{1}{2\sqrt{3}},\qquad \sqrt{1 + \theta^{2}} = -1,
\]
whence
\[
A = \dfrac{4}{\sqrt{3}} = 2{\cdot}30940;
\]
but if $c = 73^{\circ}\,54' + \epsilon$, then $\theta = \tfrac{13}{12}\epsilon$, $\theta + a = -\dfrac{1}{2\sqrt{3}}$, $\sqrt{1 + \theta^{2}} = +1$, and
\[
A = -\dfrac{4}{\sqrt{3}} = -2{\cdot}30940,
\]
viz., there is an abrupt change from $A = +\dfrac{4}{\sqrt{3}}$ to $A = -\dfrac{4}{\sqrt{3}}$; corresponding to the discontinuity of orbit already referred to. We may diminish $c$ by $180^{\circ}$, and consider the last-mentioned value, $A = -\dfrac{4}{\sqrt{3}}$, as belonging to $c = 90^{\circ} - \omega + \epsilon = -(106^{\circ}\,6' - \epsilon)$.

\paragraph{75.} Consider next that $c$ passes from $0^{\circ}$ to $-(106^{\circ}\,6' - \epsilon)$. First if $c$ is a small negative quantity $c = -\epsilon$, $\theta$ is large and positive, and $\sqrt{1 + \theta^{2}}$ having the same sign as $\theta$ (positive) is $=\theta + \dfrac{1}{2\theta}$ nearly, we have therefore $A = \dfrac{2}{3\theta}\cdot\dfrac{-1}{2\theta} = \dfrac{-1}{3\theta^{2}}$ (same as for $c = +\epsilon$). And it is easy to see that as $c$ increases negatively, $A$ is always increasing negatively, its value for $c = -90^{\circ}$ being $A = \dfrac{1 - \sqrt{13}}{3} = -{\cdot}8685$, and for $c = -106^{\circ}\,6' + \epsilon$ $A$ being $= -2{\cdot}30940$ as above. We have a diagram of $A$ (see next page).

\medskip
\noindent\textit{[p.~434]}

\medskip

\paragraph{76.} It thus appears that from $A = -\dfrac{4}{\sqrt{3}}$ to $A = -\dfrac{4}{\sqrt{3}}$, there are always to any given value of $A$ two values of $c$, or positions of the orbit-pole. In particular if $A$ be $= -1$, the curve will be a parabola; the values of $c$ lying between $0^{\circ},\, 60^{\circ}$ and between $73^{\circ}\,54',\, 90^{\circ}$ respectively.

\begin{center}
\textit{[Figure 7: schematic graph of $A$ as a function of $c$, showing the asymptotes at $c=60^{\circ}$ and at $c=-106^{\circ}\,6'$, the level $A=-1$ marked along the horizontal axis, with the abscissa scale running from $-100^{\circ}$ on the left through $0^{\circ}$ at centre to $90^{\circ}$ on the right.]}
\end{center}

To find them, writing $A = -1$, we have
\[
-3\theta - 3a = 2\theta - 2\sqrt{1 + \theta^{2}},\quad\text{that is,}\quad 5\theta + 3a = 2\sqrt{1 + \theta^{2}},
\]
or
\[
21\theta^{2} + 30a\theta + 9a^{2} - 4 = 0,
\]
that is, substituting for $a$ its value $= \dfrac{1}{2\sqrt{3}}$,
\[
21\theta^{2} + 5\sqrt{3}\theta - \tfrac{13}{4} = 0,\qquad (14\theta\sqrt{3} + 5)^{2} = 116,
\]
or
\[
\theta = \dfrac{-5 \pm \sqrt{116}}{14\sqrt{3}},
\]
giving
\[
\theta = -{\cdot}23797,\quad\text{or}\quad \theta = -{\cdot}65034,
\]
\[
\cot c = +{\cdot}93902,\quad\text{or}\quad \cot c = +{\cdot}05071,
\]
\[
c = 46^{\circ}\,48',\quad\text{or}\quad c = 87^{\circ}\,6'.
\]

\paragraph{77.} It has been seen that $c = 73^{\circ}\,54' + \epsilon$ gives $A = -\dfrac{4}{\sqrt{3}} = -2{\cdot}30940$; there will be between $0^{\circ}$ and $60^{\circ}$ another value of $c$, giving for $A$ this same value; to find this value write $A = -\dfrac{4}{\sqrt{3}}$, then we have
\[
2\theta\Bigl(1 + \dfrac{2}{\sqrt{3}}\theta\Bigr) - 4 = \dfrac{4}{\sqrt{3}}\Bigl(\sqrt{1 + \theta^{2}} - \theta\Bigr)\cdot\sqrt{3},
\]

\medskip
\noindent\textit{[p.~435]}

\medskip

that is
\[
(1 + 2\sqrt{3})\theta + 1 = \sqrt{1 + \theta^{2}},
\]
or
\[
\theta^{2}(12 + 4\sqrt{3}) + (2 + 4\sqrt{3})\theta = 0,
\]
satisfied as it should be by $\theta = 0$, and also by
\[
\theta = -\dfrac{1}{2(3 + \sqrt{3})},
\]
giving
\[
\cot c = {\cdot}76038\quad\text{or}\quad c = 52^{\circ}\,45'.
\]

\paragraph{78.} Representing the equation of the orbit by
\[
r = Ax \pm a(1 - A^{2}),
\]
we have for the point $1$,
\[
1 = A \pm a(1 - A^{2}),
\]
that is
\[
a = \dfrac{\pm 1}{1 + A},
\]
where the sign is to be taken so that $a$ shall be positive.

\paragraph{79.} With a view to the calculation of the times of passage, I calculate a series of values of $x_{2},\, y_{2},\, r_{2},\, A,\, a$, for values of $c$ at the intervals of $5^{\circ}$ and for a few intermediate values; we have $x_{3},\, y_{3},\, r_{3} = x_{2},\, y_{2},\, r_{2}$, so that these are known; so long as the orbit is an ellipse, the time of passage between the points $2$ and $3$, say $T_{23}$, may be calculated by Lambert's equation, the length of the chord $= y_{2} - y_{3}$, $= 2y_{2}$ being known without any fresh calculation. And then the times $T_{12}$ and $T_{31}$ being equal, and the sum $T_{12} + T_{23} + T_{31}$ being equal to the whole periodic time (reckoned as $= 3a^{\frac{3}{2}}$), the times $T_{12}$ and $T_{31}$ are also known. But when the orbit is a concave hyperbola there is no time $T_{23}$, and the other two times $T_{12},\, = T_{31}$ must be calculated. For the reason referred to (ante, No.~$39$) I did not use Lambert's equation, and it was less necessary to do so, by reason that, the transverse axis coinciding with the axis of $x$, the other formula could be employed without difficulty.

\paragraph{80.} The formulae for $x_{2},\, y_{2}$ adapted to logarithmic calculation are
\[
\begin{aligned}
\log(x_{2} + {\cdot}615.39) &= 11{\cdot}60174 + \log\tan(c + 16^{\circ}\,6'),\\
\log y_{2} &= n\!\cdot\!92015 + \log\sec(c + 16^{\circ}\,6'),
\end{aligned}
\]
where $y_{2}$ is always positive, but the sign of $x_{2}$ must be attended to. The values of $r_{2}$ and its inclination $\phi_{2}$ to the axis of $x$ are then to be calculated from
\[
\dfrac{y_{2}}{x_{2}} = \tan\phi_{2};\quad r_{2} = x_{2}\sec\phi_{2}\ \text{or}\ = y_{2}\csc\phi_{2},
\]
(viz.\ for $r_{2}$ it is proper to use the first or the second value, according as $x_{2}$ is greater or less than $y_{2}$).

\medskip
\noindent\textit{[p.~436]}

\medskip

We have then $e\,(= \pm A)$ and $a$ from the foregoing formulae
\[
A = \dfrac{r_{2} - 1}{x_{2} - 1},\qquad a = \dfrac{\pm 1}{1 + A},
\]
where $a,\, e$ are each of them positive.

And then for the Times
\[
T_{31} = T_{12} = \tfrac{3}{2\pi}\,a^{\frac{3}{2}}\bigl(\chi' - \chi - \sin\chi' + \sin\chi\bigr);\qquad (\log\tfrac{3}{2\pi} = \overline{1}{\cdot}67894)
\]
where
\[
\begin{aligned}
a\cos\chi &= a - r_{2} - y_{2},\\
a\cos\chi' &= a - r_{2} + y_{2},
\end{aligned}
\]
and attention is necessary in order to the selection of the proper values of the angles $\chi,\, \chi'$.

And finally
\[
T_{23} = T_{12} - \tfrac{3}{2\pi}\,(3a^{\frac{3}{2}} - T_{12}).
\]

\paragraph{81.} I subjoin a specimen; the characteristics of the logarithms are (as in the actual calculations) omitted.

\medskip
\noindent\textit{[Specimen calculation, $b=90^{\circ}$, $c=20^{\circ}$. The page presents a column-arithmetic worked example. The numerical content from the scan is reproduced below in alignment.]}

\smallskip

\begin{center}
\begin{tabular}{l@{\hspace{2em}}l}
$b = 90^{\circ}\quad c = 20^{\circ}$ & \\
$c + 16^{\circ}\,6' = 36^{\circ}\,6'$ & \\
$\log\sec\ \cdot 09259$ & $\log\tan\ \cdot 86285$\\
$\quad\ {\cdot}92015$ & $\quad\ \cdot 60174$\\
$\ \ \overline{\phantom{xxx}}$ & $\ \ \overline{\phantom{xxx}}$\\
$\quad\ {\cdot}01274$ & $\quad\ {\cdot}46459$\\
$y_{2} = 1{\cdot}0297$ & $\quad\ {\cdot}61539$\\
& $\quad\ \cdot 01274$\\
& $\quad\ \cdot 29147$\\
& $\quad\ \cdot 51044$\\
& $\therefore x_{2} = -{\cdot}32{\cdot}392$\\
& $\log = {\cdot}51044$\\
& $\quad\ {\cdot}02046$\\
& $\quad\ {\cdot}01274$\\
${\cdot}50230$ & $\quad\ {\cdot}03320$\\
$\phi_{2} = 72^{\circ}\,33'$ & $r_{2} = 1{\cdot}0794$\\
& \\
$\cdot 0794\ \log = \overline{1}{\cdot}89982$ & $\cdot 94003\ \log = \overline{1}{\cdot}97314$\\
$1{\cdot}3239\ \log = {\cdot}12185$ & $\text{comp} = {\cdot}02686$\\
& $\quad\ {\cdot}77797$\\
\end{tabular}
\end{center}

\medskip
\noindent\textit{[p.~437]}

\medskip

\begin{center}
\begin{tabular}{l@{\hspace{2em}}l}
$A = -{\cdot}059975$ & $a = 1{\cdot}0638$\\
& $\quad\ 1{\cdot}0638$\\
& $\quad\ 1{\cdot}0794$\\
$A - r_{2} = -{\cdot}0156$ & \\
$y_{2} = +1{\cdot}0297$ & \\
$\quad\ 1{\cdot}0141 = a\cos\chi'$ & \\
$-1{\cdot}0453 = a\cos\chi$ & \\
& \\
$\cdot 00608$ & $\cdot 01924$\\
$\cdot 02686$ & $\cdot 02686$\\
$\cdot 97922$ & $\cdot 99238$\\
$\chi' = 17^{\circ}\,35'$ & $\chi\,(= \text{Supp.}\ 10^{\circ}\,42') = 169^{\circ}\,18'$\quad $\chi - \chi' = 151^{\circ}\,43'$\\
\multicolumn{2}{l}{$151^{\circ}\quad 2{\cdot}63544\qquad\quad {\cdot}02686$}\\
\multicolumn{2}{l}{$\phantom{151^{\circ}\,\,}43'\quad {\cdot}01250\qquad\quad {\cdot}01343$}\\
\multicolumn{2}{l}{$-\sin\chi\quad -{\cdot}18566\qquad\ \ {\cdot}47712$}\\
\multicolumn{2}{l}{$\sin\chi'\quad {\cdot}30209\qquad\quad\ \ {\cdot}51741$}\\
\multicolumn{2}{l}{$\phantom{-\sin\chi\ \ }2{\cdot}95003\qquad\qquad\quad 3a^{\frac{3}{2}} = 3{\cdot}2916$}\\
\multicolumn{2}{l}{$\phantom{-\sin\chi\ \ \,}{\cdot}18566\qquad\qquad\qquad\qquad\ 1{\cdot}4482$}\\
\multicolumn{2}{l}{$\phantom{-\sin\chi\ \ }2{\cdot}76437\quad\log = {\cdot}44160\qquad 1{\cdot}8434$}\\
\multicolumn{2}{l}{$\phantom{-\sin\chi\ \ \ \ \ \ \ }{\cdot}02686\qquad\quad T_{12} = T_{31} = {\cdot}9217$}\\
\multicolumn{2}{l}{$\phantom{-\sin\chi\ \ \ \ \ \ \ }{\cdot}01343$}\\
\multicolumn{2}{l}{$\phantom{-\sin\chi\ \ \ \ \ \ \ }{\cdot}67894$}\\
\multicolumn{2}{l}{$T_{23} = 1{\cdot}4482\qquad {\cdot}16083$}\\
\end{tabular}
\end{center}

\paragraph{82.} For the Time in a hyperbola, we have
\[
T_{31} = T_{12} = \dfrac{3}{2\pi}\,a^{\frac{3}{2}}\bigl\{e\tan u - h\!\cdot\!l.\tan(45^{\circ} + \tfrac{1}{2}u)\bigr\},
\]
where
\[
\tan u = \dfrac{y_{2}}{a\sqrt{e^{2} - 1}}.
\]

\medskip
\noindent\textit{[p.~438]}

\medskip

Taking as a specimen the case $c = 75^{\circ}$, we have here
\[
\begin{aligned}
a &= {\cdot}9004 & e &= 2{\cdot}1106 & y_{2} &= 43{\cdot}341\\
\log &= \overline{1}{\cdot}95444 & \log &= {\cdot}32441 & \log &= 1{\cdot}63690
\end{aligned}
\]
\[
\begin{aligned}
a(e^{2} - 1) &= 3{\cdot}1106\\
\log &= {\cdot}49284
\end{aligned}
\]
and then the calculation is
\begin{center}
\begin{tabular}{l@{\hspace{2em}}l}
$\log a = \overline{1}{\cdot}95444$ & \\
$\log a(e^{2} - 1) = {\cdot}49284$ & \\
$\ \ \overline{\phantom{xxx}}\ {\cdot}44728$ & \\
$\log a\sqrt{e^{2} - 1} = {\cdot}22364$ & \\
$\log y_{2} = 1{\cdot}63690$ & \\
$\log\tan u = 1{\cdot}41326$ & $u = 87^{\circ}\,47'$\\
& $h\!\cdot\!l\tan(45^{\circ} + \tfrac{1}{2}u) = 3{\cdot}95140$\\
$\ \ {\cdot}32441$ & $e\tan u = 54{\cdot}660$\\
$\ \ \overline{\phantom{xxx}}\ {\cdot}73767$ & $\quad\ 3{\cdot}951$\\
& $\quad\ 50{\cdot}709$\\
& $\log = {\cdot}70508$\\
& $\quad\ {\cdot}95444$\\
& $\quad\ {\cdot}47722$\\
& $\quad\ {\cdot}67894$\\
& \\
& $\quad\ 3{\cdot}1568$\\
& $T_{12} = T_{31} = 20{\cdot}686$\\
\end{tabular}
\end{center}

\paragraph{83.} In the case of the parabola $p = 1$, and the expression for the Times is
\[
T_{12} = T_{31} = \dfrac{1}{6\sqrt{\pi}}\Bigl\{(\rho + \rho' + \gamma)^{\frac{3}{2}} - (\rho + \gamma)^{\frac{3}{2}}\Bigr\},
\]
where for
\[
\begin{aligned}
c &= 46^{\circ}\,48'\\
c &= 87^{\circ}\,6'
\end{aligned}
\]
we have
\[
\begin{aligned}
T_{12} = T_{31} &= {\cdot}787,\\
T_{12} = T_{31} &= 2{\cdot}588.
\end{aligned}
\]

\medskip
\noindent\textit{[p.~439]}

\medskip

\begin{center}
\textit{Planogram No.~$1$, first part, $b = 90^{\circ}$.}

\smallskip

\begin{tabular}{l@{\hspace{0.4em}}r@{\hspace{0.4em}}r@{\hspace{0.4em}}r@{\hspace{0.4em}}r@{\hspace{0.4em}}r@{\hspace{0.4em}}r@{\hspace{0.4em}}r@{\hspace{0.4em}}r@{\hspace{0.4em}}r@{\hspace{0.4em}}r@{\hspace{0.4em}}r}
\hline
& $c$ & $x_{2}$ & $y_{2}$ & $\phi_{2}$ & $r_{2}$ & $A$ & $a$ & $T_{12}$ & $T_{23}$ & $T_{31}$\\
\hline
Circle & $0^{\circ}$ & $-{\cdot}500$ & $+{\cdot}866$ & $60^{\circ}$ & $1{\cdot}000$ & $0$ & $1{\cdot}000$ & $1{\cdot}000$ & $1{\cdot}000$ & $1{\cdot}000$\\
 & $5$ & ${\cdot}461$ & ${\cdot}892$ & $62^{\circ}\,39'$ & $1{\cdot}004$ & $-{\cdot}003$ & $1{\cdot}003$ & & & \\
 & $10$ & ${\cdot}420$ & ${\cdot}927$ & $65\,38$ & $1{\cdot}017$ & $-{\cdot}012$ & $1{\cdot}012$ & ${\cdot}960$ & $1{\cdot}135$ & ${\cdot}960$\\
 & $15$ & ${\cdot}374$ & ${\cdot}972$ & $68\,56$ & $1{\cdot}041$ & $-{\cdot}030$ & $1{\cdot}031$ & & & \\
Ellipses & $20$ & ${\cdot}324$ & $1{\cdot}030$ & $72\,33$ & $1{\cdot}079$ & $-{\cdot}060$ & $1{\cdot}064$ & ${\cdot}922$ & $1{\cdot}448$ & ${\cdot}922$\\
 & $25$ & ${\cdot}267$ & $1{\cdot}104$ & $76\,26$ & $1{\cdot}136$ & $-{\cdot}107$ & $1{\cdot}120$ & ${\cdot}887$ & $2{\cdot}275$ & ${\cdot}887$\\
 & $30$ & ${\cdot}200$ & $1{\cdot}200$ & $80\,32$ & $1{\cdot}216$ & $-{\cdot}180$ & $1{\cdot}220$ & & & \\
 & $35$ & ${\cdot}120$ & $1{\cdot}325$ & $84\,49$ & $1{\cdot}330$ & $-{\cdot}295$ & $1{\cdot}418$ & & & \\
 & $40$ & $-{\cdot}021$ & $1{\cdot}492$ & $90\,48$ & $1{\cdot}492$ & $-{\cdot}482$ & $1{\cdot}931$ & ${\cdot}838$ & $6{\cdot}371$ & ${\cdot}838$\\
 & $40^{\circ}\,54'$ & ${\cdot}000$ & $1{\cdot}515$ & $90$ & $1{\cdot}515$ & $-{\cdot}515$ & $2{\cdot}061$ & & & \\
 & $45$ & $+{\cdot}109$ & $1{\cdot}722$ & $93\,37$ & $1{\cdot}725$ & $-{\cdot}814$ & $5{\cdot}362$ & & & \\
Parab. & $46^{\circ}\,48'$ & ${\cdot}166$ & $1{\cdot}826$ & $95\,11$ & $1{\cdot}834$ & $-1{\cdot}000$ & $\infty$ & ${\cdot}787$ & $\infty$ & ${\cdot}787$\\
 & $50$ & ${\cdot}287$ & $2{\cdot}054$ & $97\,57$ & $2{\cdot}074$ & $-1{\cdot}505$ & $1{\cdot}981$ & ${\cdot}750$ & $\sim$ & ${\cdot}750$\\
Hyperbs.\ \{ & $52^{\circ}\,45'$ & ${\cdot}418$ & $2{\cdot}306$ & $100\,16$ & $2{\cdot}344$ & $-2{\cdot}309$ & ${\cdot}764$ & ${\cdot}628$ & $\sim$ & ${\cdot}628$\\
 & $55$ & ${\cdot}552$ & $2{\cdot}569$ & $102\,8$ & $2{\cdot}627$ & $-3{\cdot}632$ & ${\cdot}380$ & & & \\
Line & $60$ & $1{\cdot}000$ & $3{\cdot}464$ & $106\,6$ & $3{\cdot}605$ & $\mp\infty$ & ${\cdot}000$ & ${\cdot}000$ & $\sim$ & ${\cdot}000$\\
 & $65$ & $1{\cdot}937$ & $5{\cdot}378$ & $109\,48$ & $5{\cdot}716$ & $+5{\cdot}032$ & ${\cdot}166$ & & Convex orbit. & \\
Convex & $70$ & $5{\cdot}248$ & $12{\cdot}233$ & $113\,13$ & $13{\cdot}311$ & $+2{\cdot}898$ & ${\cdot}257$ & & & \\
 & $73^{\circ}\,54'-\epsilon$ & $+\infty$ & $\infty$ & $115\,39$ & $\infty$ & $+2{\cdot}309$ & ${\cdot}302$ & & & \\
 & $73^{\circ}\,54'+\epsilon$ & $-\infty$ & $\infty$ & $64\,21$ & $\infty$ & $-2{\cdot}309$ & ${\cdot}764$ & $\infty$ & $\sim$ & $\infty$\\
 & $76$ & $-21{\cdot}432$ & $43{\cdot}341$ & $63\,42$ & $48{\cdot}346$ & $-2{\cdot}111$ & ${\cdot}900$ & $20{\cdot}68$ & $\sim$ & $20{\cdot}68$\\
Hyperbs. & $80$ & $-4{\cdot}356$ & $7{\cdot}830$ & $60\,55$ & $8{\cdot}960$ & $-1{\cdot}486$ & $2{\cdot}056$ & $3{\cdot}856$ & $\sim$ & $3{\cdot}856$\\
 & $85$ & $-2{\cdot}653$ & $4{\cdot}322$ & $58\,27$ & $5{\cdot}072$ & $-1{\cdot}115$ & $8{\cdot}718$ & $2{\cdot}912$ & $\sim$ & $2{\cdot}912$\\
Parab. & $87^{\circ}\,6'$ & $-2{\cdot}320$ & $3{\cdot}644$ & $57\,31$ & $4{\cdot}320$ & $-1{\cdot}000$ & $\infty$ & $2{\cdot}588$ & $\infty$ & $2{\cdot}588$\\
Ellipse & $90$ & $-2{\cdot}000$ & $3{\cdot}000$ & $56\,18$ & $3{\cdot}606$ & $-{\cdot}869$ & $7{\cdot}622$ & $2{\cdot}255$ & $58{\cdot}62$ & $2{\cdot}255$\\
\hline
\end{tabular}

\smallskip

\noindent\small The mark $\sim$ in the $T_{23}$ column shows that there is no Time $T_{23}$.
\end{center}

\medskip
\noindent\textit{[p.~440]}

\medskip

\begin{center}
\textit{Planogram No.~$1$, second part, $b = 270^{\circ}$.}

\smallskip

\begin{tabular}{l@{\hspace{0.4em}}r@{\hspace{0.4em}}r@{\hspace{0.4em}}r@{\hspace{0.4em}}r@{\hspace{0.4em}}r@{\hspace{0.4em}}r@{\hspace{0.4em}}r@{\hspace{0.4em}}r@{\hspace{0.4em}}r@{\hspace{0.4em}}r@{\hspace{0.4em}}r}
\hline
& $c$ & $x_{2}$ & $y_{2}$ & $\phi_{2}$ & $r_{2}$ & $A$ & $a$ & $T_{12}$ & $T_{23}$ & $T_{31}$\\
\hline
Circ. & $0^{\circ}$ & $-{\cdot}500$ & $+{\cdot}866$ & $60^{\circ}$ & $1{\cdot}000$ & $0$ & $1{\cdot}000$ & $1{\cdot}000$ & $1{\cdot}000$ & $1{\cdot}000$\\
 & $5$ & ${\cdot}537$ & ${\cdot}848$ & $57\,39$ & $1{\cdot}003$ & $-{\cdot}002$ & $1{\cdot}002$ & & & \\
 & $10$ & ${\cdot}573$ & ${\cdot}837$ & $55\,37$ & $1{\cdot}014$ & $-{\cdot}009$ & $1{\cdot}009$ & $1{\cdot}044$ & ${\cdot}951$ & $1{\cdot}044$\\
 & $15$ & ${\cdot}608$ & ${\cdot}832$ & $53\,52$ & $1{\cdot}030$ & $-{\cdot}019$ & $1{\cdot}019$ & & & \\
 & $20$ & ${\cdot}643$ & ${\cdot}834$ & $52\,23$ & $1{\cdot}053$ & $-{\cdot}032$ & $1{\cdot}033$ & $1{\cdot}091$ & ${\cdot}969$ & $1{\cdot}091$\\
 & $25$ & ${\cdot}678$ & ${\cdot}842$ & $51\,10$ & $1{\cdot}081$ & $-{\cdot}048$ & $1{\cdot}051$ & & & \\
 & $30$ & ${\cdot}714$ & ${\cdot}857$ & $50\,11$ & $1{\cdot}116$ & $-{\cdot}068$ & $1{\cdot}073$ & $1{\cdot}145$ & $1{\cdot}043$ & $1{\cdot}145$\\
All Ellipses & $35$ & ${\cdot}752$ & ${\cdot}879$ & $49\,27$ & $1{\cdot}157$ & $-{\cdot}090$ & $1{\cdot}098$ & & & \\
 & $40$ & ${\cdot}793$ & ${\cdot}910$ & $48\,51$ & $1{\cdot}207$ & $-{\cdot}115$ & $1{\cdot}130$ & $1{\cdot}207$ & $1{\cdot}192$ & $1{\cdot}207$\\
 & $45$ & ${\cdot}836$ & ${\cdot}950$ & $48\,40$ & $1{\cdot}266$ & $-{\cdot}145$ & $1{\cdot}169$ & & & \\
 & $50$ & ${\cdot}884$ & $1{\cdot}002$ & $48\,36$ & $1{\cdot}336$ & $-{\cdot}179$ & $1{\cdot}217$ & $1{\cdot}283$ & $1{\cdot}464$ & $1{\cdot}283$\\
 & $55$ & ${\cdot}938$ & $1{\cdot}069$ & $48\,45$ & $1{\cdot}442$ & $-{\cdot}218$ & $1{\cdot}278$ & & & \\
 & $60$ & $1{\cdot}000$ & $1{\cdot}154$ & $49\,7$ & $1{\cdot}527$ & $-{\cdot}264$ & $1{\cdot}358$ & $1{\cdot}377$ & $1{\cdot}983$ & $1{\cdot}377$\\
 & $65$ & $1{\cdot}074$ & $1{\cdot}266$ & $49\,42$ & $1{\cdot}660$ & $-{\cdot}318$ & $1{\cdot}466$ & & & \\
 & $70$ & $1{\cdot}164$ & $1{\cdot}412$ & $50\,31$ & $1{\cdot}830$ & $-{\cdot}383$ & $1{\cdot}622$ & $1{\cdot}506$ & $3{\cdot}036$ & $1{\cdot}506$\\
 & $75$ & $1{\cdot}280$ & $1{\cdot}611$ & $51\,35$ & $2{\cdot}056$ & $-{\cdot}464$ & $1{\cdot}864$ & & & \\
 & $80$ & $1{\cdot}431$ & $1{\cdot}891$ & $52\,53$ & $2{\cdot}372$ & $-{\cdot}564$ & $2{\cdot}295$ & $1{\cdot}771$ & $6{\cdot}888$ & $1{\cdot}771$\\
 & $85$ & $1{\cdot}651$ & $2{\cdot}311$ & $54\,28$ & $2{\cdot}840$ & $-{\cdot}694$ & $3{\cdot}269$ & & & \\
 & $90$ & $-2{\cdot}000$ & $+3{\cdot}000$ & $56\,18$ & $3{\cdot}606$ & $-{\cdot}869$ & $7{\cdot}622$ & $2{\cdot}255$ & $58{\cdot}62$ & $2{\cdot}255$\\
\hline
\end{tabular}
\end{center}

\medskip
\noindent\textit{[p.~441]}

\medskip

\begin{center}
\textit{Article Nos.~$84$ to $94$. Planogram No.~$2$, the Meridian $0^{\circ}$--$180^{\circ}$ (see Plate III.).}
\end{center}

\paragraph{84.} The orbit-plane here rotates about an axis in the plane of the ecliptic at right angles to $S1$ (Fig.~$6$). The entire meridian is given by $b = 0^{\circ}$, $c = 0^{\circ}$ to $90^{\circ}$, and $b = 180^{\circ}$, $c = 0^{\circ}$ to $90^{\circ}$, but it is sufficient to consider one of these half meridians, say the latter of them, as the series of values is the same for each of them, with only an interchange of the points $2,\, 3$. I write therefore, $b = 180^{\circ}$, so that we have
\[
\begin{aligned}
\alpha,\ \beta,\ \gamma &= 0,\ 1,\ 0,\\
\alpha',\ \beta',\ \gamma' &= -\cos c,\ 0,\ -\sin c,\\
\alpha'',\, \beta'',\, \gamma'' &= -\sin c,\ 0,\ \cos c,
\end{aligned}
\]
consequently
\[
\begin{aligned}
x_{1}'\,:\,y_{1}'\,:\,1 &= \sin c\ :\ -\sqrt{3}\ :\ \sqrt{3}\cos c,\\
x_{2}'\,:\,y_{2}'\,:\,1 &= -1 - 3\cos c - 2\sin c\ :\ -\sqrt{3}\ :\ -\sqrt{3}\sin c - 2\sqrt{3}\cos c,\\
x_{3}'\,:\,y_{3}'\,:\,1 &= -1 - 3\cos c - 2\sin c\ :\ -\sqrt{3}\ :\ \sqrt{3}\sin c - 2\sqrt{3}\cos c,
\end{aligned}
\]
and moreover $x = y',\ y = -x'$; so that, introducing into the formulae $(x_{1},\, y_{1})$, \&c., in place of the $(x_{1}',\, y_{1}')$, \&c., we have
\[
x_{1} = \sec c,\qquad y_{1} = -\tfrac{1}{\sqrt{3}}\tan c,
\]
\[
x_{2} = \dfrac{-1}{\sin c - 2\cos c}\,(\sin c + 2\cos c) + \tfrac{1}{\sqrt{3}}\cdot\dfrac{1 + 2\sin c - 3\cos c}{\sqrt{3}\sin c - 2\cos c},
\]
\[
x_{3} = \dfrac{-1}{\sin c - 2\cos c}\,(\sin c + 2\cos c) - \tfrac{1}{\sqrt{3}}\cdot\dfrac{1 - 2\sin c - 3\cos c}{\sqrt{3}\sin c - 2\cos c},
\]
which, putting
\[
\cos\delta = \tfrac{2}{\sqrt{5}},\quad \sin\delta = \tfrac{1}{\sqrt{5}},\quad \tan\delta = \tfrac{1}{2},\quad \delta = 26^{\circ}\,34',
\]
become
\[
\begin{aligned}
x_{1} &= \sec c,\qquad\qquad\quad\, y_{1} = -\tfrac{1}{\sqrt{3}}\tan c,\\
x_{2} &= -\tfrac{1}{\sqrt{5}}\sec(c - \delta),\quad y_{2} = \tfrac{1}{\sqrt{3}}\bigl\{\tfrac{2}{5} + \tfrac{1}{2}\tan(c - \delta)\bigr\},\\
x_{3} &= -\tfrac{1}{\sqrt{5}}\sec(c + \delta),\quad y_{3} = \tfrac{1}{\sqrt{3}}\bigl\{-\tfrac{2}{5} + \tfrac{1}{2}\tan(c + \delta)\bigr\}
\end{aligned}
\]
so that the guide-hyperbolas are (angle of asymptotes $= 30^{\circ}$):
\[
\begin{aligned}
x_{1}^{2} - 3y_{1}^{2} &= 1,\\
x_{2}^{2} &= 15y_{2}^{2} - 16\sqrt{3}\,y_{2} + 13,\qquad \tan^{-1}\tfrac{1}{\sqrt{15}} = 14^{\circ}\,28'\\
x_{3}^{2} &= 15y_{3}^{2} + 16\sqrt{3}\,y_{3} + 13,\qquad
\end{aligned}
\]

\medskip
\noindent\textit{[p.~442]}

\medskip

It is easy to verify that

Hyperbola $2$ passes through $x_{2} = -\tfrac{1}{2}$, $y_{2} = +\tfrac{1}{2}\sqrt{3}$, and touches there circle $x^{2} + y^{2} = 1$,

\quad\,\,$3$\quad\quad\quad\,\,$x_{3} = -\tfrac{1}{2}$, $y_{3} = -\tfrac{1}{2}\sqrt{3}$\quad\quad\quad\quad\quad\quad\quad\quad ,,

and we thus have the figure in the Plate.

\paragraph{85.} The figure shows the motion of the points $1$, $2$, $3$, along their respective hyperbolas, viz.\ $c = 0^{\circ}$ to $90^{\circ}$, the point $1$ moves from contact with the circle, along a half branch to infinity: $2$ moves from contact along a small portion of the half branch; $3$ moves from contact, along the half branch to infinity for $c = \tan^{-1}2 = 63^{\circ}\,26'$, and then reappearing at the opposite infinity, as $c$ increases to $90^{\circ}$ describes a portion of the opposite half branch.

\paragraph{86.} For $c = 0$, the orbit is the circle; as $c$ increases the orbit becomes elliptic, then parabolic, $c = 51^{\circ}$, and afterwards hyperbolic (concave); until for $c = 60^{\circ}$, the three points are on the horizontal line of the figure, and the orbit is this right line; it is to be noticed that the arrangement of the points on these orbits is $1,\, 2,\, 3$; so that for the parabola, $T_{23}$ is $= \infty$, and for the hyperbolas and right line $T_{23}$ does not exist.

\paragraph{87.} For $c < 60^{\circ}$ until $c = 63^{\circ}\,26'$ the orbit is a convex hyperbola, the arrangement of the points being still $1,\, 2,\, 3$: say for $c = 63^{\circ}\,26' - \epsilon$, the orbit is the convex hyperbola $\Omega$. At $c = 63^{\circ}\,26'$ there is an abrupt change of orbit; say for $c = 63^{\circ}\,26' + \epsilon$ the orbit is a concave hyperbola $\Omega_{1}$; and for $c = 65^{\circ}\,52'$ the orbit is a parabola; the arrangement of the points on these orbits is $2,\, 1,\, 3$; so that for the hyperbolas $T_{23}$ does not exist, and $T_{21}$ is $= \infty$ for the parabola. Observe also that for the hyperbola $\Omega_{1}$, the point $3$ is at infinity, or we have $T_{31} = \infty$. As $c$ continues to increase, the orbit becomes an ellipse, the eccentricity having a minimum value $= {\cdot}628$ (about), for $c = 69^{\circ}$ (about). For $c = 89^{\circ}\,20'$ the orbit is again a parabola, and then until $c = 90^{\circ}$ it is a hyperbola; the order of the points on the last-mentioned parabola and hyperbolas being $1,\, 3,\, 2$; so that for the parabola $T_{32}$ is $= \infty$, and for the hyperbolas $T_{32}$ does not exist. In the hyperbola for $c = 90^{\circ}$, say the hyperbola $\Omega'$, the point $1$ is at infinity, or we have $T_{12} = \infty$. The foregoing results, obtained (except as to the numerical values) by consideration of the figure, will be confirmed by means of the calculated values of $e$.

\paragraph{88.} The equation of the orbit may be written
\[
\resizebox{\textwidth}{!}{$\displaystyle
\left|\begin{array}{cccc}
\dfrac{r}{\sqrt{3}} & \dfrac{x}{\sqrt{3}} & y & \dfrac{1}{\sqrt{3}}\\[6pt]
r_{1}\cos c & ,\quad 1, & \sin c & \cos c\\[4pt]
+r_{2}(\sin c + 2\cos c), & -1, & +2\sin c - 3\cos c, & \sin c + 2\cos c\\[2pt]
& & & \sin c - 2\cos c\\
-r_{3}(\sin c + 2\cos c), & 1, & +2\sin c - 3\cos c, & \\
\end{array}\right| = 0,
$}
\]

\medskip
\noindent\textit{[p.~443]}

\medskip

or developing, this is
\[
\resizebox{\textwidth}{!}{$\displaystyle
\begin{aligned}
&\tfrac{r}{\sqrt{3}}\,6(\sin^{2}c - 3\cos^{2}c)\\
&-\tfrac{x}{\sqrt{3}}\Bigl\{\,4r_{1}(\sin^{2}c - 3\cos^{2}c)\cos c\\
&\hspace{1.4cm} - r_{2}(\sin c + 2\cos c)(\sin^{2}c - 3\cos^{2}c)\\
&\hspace{1.4cm} + r_{3}(\sin c - 2\cos c)(\sin^{2}c - 3\cos^{2}c)\Bigr\}\\
&+y\,\bigl\{\,-2r_{1}\sin c\cos c\\
&\hspace{1.4cm} + r_{2}(-\sin^{2}c + \sin c\cos c + 6\cos^{2}c)\\
&\hspace{1.4cm} + r_{3}(\,\sin^{2}c + \sin c\cos c - 6\cos^{2}c)\bigr\}\\
&-\tfrac{1}{\sqrt{3}}\Bigl\{\,r_{1}\cdot 6\cos^{2}c\\
&\hspace{1.4cm} + 3r_{2}(\sin^{2}c + \sin c\cos c + 2\cos^{2}c)\\
&\hspace{1.4cm} + 3r_{3}(\sin^{2}c - \sin c\cos c - 2\cos^{2}c)\Bigr\} = 0\,;
\end{aligned}$}
\]
(observe that the orbit will be a right line if $\sin^{2}c - 3\cos^{2}c = 0$, that is if $c = 60^{\circ}$, which is right, since $60^{\circ}$ is the angular radius of the regulator circle).

\paragraph{89.} Putting in the equation $\tan c = \lambda$, and therefore $\cos c = \dfrac{1}{\sqrt{1 + \lambda^{2}}}$, the equation becomes
\[
\resizebox{\textwidth}{!}{$\displaystyle
r = \dfrac{1}{6\sqrt{1 + \lambda^{2}}}\bigl(4r_{1} - (\lambda + 2)r_{2} + (\lambda - 2)r_{3}\bigr)x
+ \dfrac{1}{2\sqrt{3}(\lambda^{2} - 3)}\bigl(2\lambda r_{1} + (\lambda - 3)(\lambda - 2)r_{2} - (\lambda + 3)(\lambda + 2)r_{3}\bigr)y
+ \dfrac{1}{2(\lambda^{2} - 3)}\bigl(-2r_{1} + (\lambda - 1)(\lambda + 2)r_{2} + (\lambda + 1)(\lambda - 2)r_{3}\bigr).
$}
\]
We have
\[
x_{1} = \sqrt{1 + \lambda^{2}},\quad x_{2} = \dfrac{-\sqrt{1 + \lambda^{2}}}{\lambda + 2},\quad x_{3} = \dfrac{-\sqrt{1 + \lambda^{2}}}{\lambda - 2},
\]
\[
y_{1} = \tfrac{1}{\sqrt{3}}\,\lambda,\quad y_{2} = \tfrac{1}{\sqrt{3}}\cdot\dfrac{2\lambda + 3}{\lambda + 2},\quad y_{3} = \tfrac{1}{\sqrt{3}}\cdot\dfrac{2\lambda - 3}{\lambda - 2},
\]
and thence, writing for shortness
\[
\begin{aligned}
R_{1} &= \sqrt{1 + \tfrac{4}{3}\lambda^{2}},\\
R_{2} &= \tfrac{1}{\sqrt{3}}\sqrt{7\lambda^{2} + 12\lambda + 12},\\
R_{3} &= \tfrac{1}{\sqrt{3}}\sqrt{7\lambda^{2} - 12\lambda + 12},
\end{aligned}
\]

\medskip
\noindent\textit{[p.~444]}

\medskip

we have
\[
r_{1} = R_{1},\quad r_{2}(\lambda + 2) = R_{2},\quad r_{3}(\lambda - 2) = R_{3},
\]
where $r_{1},\, r_{2},\, r_{3}$ are positive, and the signs of $R_{1},\, R_{2},\, R_{3}$ must be determined accordingly; viz., $R_{1}$ is always positive, and ($c = 0^{\circ}$ to $c = 90^{\circ}$, as here supposed) $R_{2}$ is also positive but $R_{3}$ has the same sign as $\lambda - 2$; viz., $c = 0^{\circ}$ to $c = 63^{\circ}\,26'$, $R_{3}$ is negative; and $c = 63^{\circ}\,26'$ to $c = 90^{\circ}$, $R_{3}$ is positive. It is to be observed that this position, $c = \tan^{-1}2 = 63^{\circ}\,26'$, of the pole is the intersection of the meridian $b = 180^{\circ}$ by a separator circle, and corresponds to an intersection at infinity on the ray $3$.

\paragraph{90.} Substituting the foregoing values of $r_{1},\, r_{2},\, r_{3}$, the equation of the orbit becomes
\[
\resizebox{\textwidth}{!}{$\displaystyle
r = \dfrac{1}{6\sqrt{1 + \lambda^{2}}}\bigl(4R_{1} - R_{2} + R_{3}\bigr)x
+ \dfrac{1}{2\sqrt{3}(\lambda^{2} - 3)}\bigl\{\lambda(2R_{1} + R_{2} - R_{3}) - 3(R_{2} + R_{3})\bigr\}\,y
+ \dfrac{1}{2(\lambda^{2} - 3)}\bigl[\lambda(R_{2} + R_{3}) - 2R_{1} - R_{2} + R_{3}\bigr],
$}
\]
where $\lambda = \tan c$; and the equation of the orbit may thence be calculated for any given value of $c$.

\paragraph{91.} The analytical expression for the eccentricity is
\[
e = \sqrt{A^{2} + B^{2}},
\]
\begin{center}
\textit{[Figure 8: graph of the eccentricity $e$ versus $c$, showing the values $e=0$ at $c=0^{\circ}$, the parabolic levels $e=1$ marked, the vertical asymptote near $c=63^{\circ}\,26'$ between the convex and concave hyperbola branches, the minimum near $c=69^{\circ}$ at $e\approx{\cdot}628$, and the second parabolic level near $c=89^{\circ}\,20'$. Abscissa runs $0^{\circ}$ to $90^{\circ}$ with $1$-unit ordinate marks.]}
\end{center}
where, as above,
\[
\begin{aligned}
A &= \dfrac{1}{6\sqrt{1 + \lambda^{2}}}\bigl(4R_{1} - R_{2} + R_{3}\bigr),\\
B &= \dfrac{1}{2\sqrt{3}(\lambda^{2} - 3)}\bigl\{\lambda(2R_{1} + R_{2} - R_{3}) - 3(R_{2} + R_{3})\bigr\};
\end{aligned}
\]

\medskip
\noindent\textit{[p.~445]}

\medskip

but this expression is too complicated to allow of an analytical discussion of the series of values of $e$ (such as was given for $A$, $\epsilon$, in planogram No.~$1$). The numerical calculation gives the results mentioned ante No.~$87$, viz., $c = 0$, $e = 0$; $c = 51^{\circ}$, $e = 1$; $c = 60^{\circ}$, $e = \infty$; $c = 63^{\circ}\,26' - \epsilon$, $e = 4{\cdot}912$; $c = 63^{\circ}\,26' + \epsilon$, $e = 1{\cdot}853$; $c = 69^{\circ}$, $e = {\cdot}628$ (min.); $c = 89^{\circ}\,20'$ (viz. $\lambda = 86{\cdot}176$), $e = 1$; $c = 90^{\circ}$, $e = 1{\cdot}018$; values which are exhibited in the diagram in the preceding page.

\paragraph{92.} It may be further remarked, in reference to the formula $r = Ax + By + C$, that for $c = 60^{\circ}$, that is $\lambda = \sqrt{3}$, we have $A$ and $C$ each infinite, but $B$ finite, and of opposite signs; viz., the equation becomes $r = -{\cdot}2242\,x + \infty(y - 1)$, that is $y = 1$, orbit a right line as above.

The abrupt change at $c = 63^{\circ}\,26'$, $\lambda = 2$, arises from the change of sign of $R_{3}$; viz., $c = 63^{\circ}\,26' - \epsilon$, $R_{3} = -2{\cdot}309$, but $c = 63^{\circ}\,26' + \epsilon$, $R_{3} = +2{\cdot}309$; the two orbits are
\[
\begin{aligned}
c = 63^{\circ}\,26' - \epsilon, &\quad r = +{\cdot}234\,x + 4{\cdot}906\,y - 3{\cdot}671, &\quad e = 4{\cdot}912,\quad a = {\cdot}159,\\
c = 63^{\circ}\,26' + \epsilon, &\quad r = {\cdot}578\,x - 1{\cdot}761\,y + 3{\cdot}257, &\quad e = 1{\cdot}853,\quad a = 1{\cdot}338.
\end{aligned}
\]

For $c = 90^{\circ}$ the equation is
\[
r = {\cdot}770\,x + {\cdot}666\,y + 1{\cdot}527
\]
and therefore $e = \sqrt{\tfrac{}{}} = 1{\cdot}018$ as above; $a = \tfrac{9}{2}\sqrt{21} = 41{\cdot}243$.

It is to be added that for $c$ nearly $= 90^{\circ}$, or $\lambda$ very large, we have
\[
R_{1} = \tfrac{2}{\sqrt{3}}\lambda,\quad R_{2} = \sqrt{\tfrac{7}{3}}\lambda + 2\sqrt{3},\quad R_{3} = \sqrt{\tfrac{7}{3}}\lambda - 2\sqrt{3},
\]
and thence
\[
\begin{aligned}
A &= \tfrac{4}{3\sqrt{3}} - \tfrac{2}{\sqrt{21}}\cdot\dfrac{1}{\lambda} = {\cdot}770 - {\cdot}430\dfrac{1}{\lambda},\\
B &= \tfrac{2}{\sqrt{7}}\dfrac{1}{\lambda} - \tfrac{6}{\sqrt{3}}\dfrac{1}{\lambda^{2}} = {\cdot}666 - 1{\cdot}890\dfrac{1}{\lambda^{2}},\\
C &= \sqrt{\tfrac{7}{3}} - \tfrac{4}{\sqrt{3}}\dfrac{1}{\lambda} = 1{\cdot}527 - 1{\cdot}555\dfrac{1}{\lambda}.
\end{aligned}
\]
It was, in fact, by means of these expressions that the value $\lambda = 86{\cdot}176$ ($c = 89^{\circ}\,20'$) corresponding to the last-mentioned parabolic orbit was obtained.

\medskip
\noindent\textit{[p.~446]}

\medskip

\paragraph{93.} For the calculation of the table we have
\[
\begin{aligned}
\log x_{1} &= \overline{10} + \log\sec c,\\
\log y_{1} &= \overline{10}{\cdot}76144 + \log\tan c,\\
\log x_{2} &= \overline{10}{\cdot}65052 + \log\sec(c - 26^{\circ}\,34'),\\
\log(y_{2} - {\cdot}92376) &= \overline{10}{\cdot}06247 + \log\tan(c - 26^{\circ}\,34'),\\
\log x_{3} &= \overline{10}{\cdot}65052 + \log\sec(c + 26^{\circ}\,34'),\\
\log(y_{3} + {\cdot}92376) &= \overline{10}{\cdot}06247 + \log\sec(c + 26^{\circ}\,34'),
\end{aligned}
\]
the values of $r_{1},\, r_{2},\, r_{3}$, are then calculated from
\[
x_{1} = r_{1}\cos\phi_{1},\quad y_{1} = r_{1}\sin\phi_{1},
\]
or say
\[
\dfrac{y_{1}}{x_{1}} = \tan\phi_{1},\quad r_{1} = x_{1}\sec\phi_{1},\ \&c.
\]
and those of the chords $\gamma_{12},\, \gamma_{23},\, \gamma_{31}$, from
\[
x_{1} - x_{2} = \gamma_{12}\cos\theta_{12},\quad y_{1} - y_{2} = \gamma_{12}\sin\theta_{12},
\]
or say
\[
\dfrac{y_{1} - y_{2}}{x_{1} - x_{2}} = \tan\theta_{12},\quad \gamma_{12} = (x_{1} - x_{2})\sec\theta_{12}.
\]

We have then to find the equation of the orbit $r = Ax + By + C$; this might be done by substituting in the determinant expression the numerical values of $x_{1},\, y_{1},\, r_{1},\, x_{2},\, y_{2},\, r_{2}$, and $x_{3},\, y_{3},\, r_{3}$, so calculating the result, but I have preferred to employ the formula of No.~$90$, using only the calculated values of $r_{1},\, r_{2},\, r_{3}$; viz.\ we have
\[
r_{1} = R_{1},\quad r_{2}(\lambda + 2) = R_{2},\quad r_{3}(\lambda - 2) = R_{3}.
\]
which gives the values of $R_{1},\, R_{2},\, R_{3}$. And then we have $e,\, \varpi,\, a$, from the equations
\[
A = e\cos\varpi,\quad B = e\sin\varpi,\quad a = \dfrac{\pm\,C}{1 - e^{2}},
\]
$e$ and $a$ being each regarded as positive. The times in the elliptic, and parabolic orbits are then calculated from Lambert's equation, as explained in regard to Planogram No.~$1$, but for the hyperbolic orbits, the other formulae were made use of.

\medskip
\noindent\textit{[p.~447]}

\medskip

\paragraph{94.} I annex a specimen; the characteristics of the logarithms are omitted.

\smallskip

\noindent\textit{[$c = 20^{\circ}$. The columnar specimen calculation gives, in summary, the following intermediate and final results.]}

\smallskip

\begin{center}
\begin{tabular}{lll}
$x_{1} = 1{\cdot}06418$ & $x_{2} = -{\cdot}45017$ & $x_{3} = -{\cdot}65049$\\
$y_{1} = {\cdot}21014$ & $y_{2} = +{\cdot}91038$ & $y_{3} = {\cdot}80171$\\
$\phi_{1} = 11^{\circ}\,10'$ & $\phi_{2}\,(= 63^{\circ}\,41') = 116^{\circ}\,19'$ & $\phi_{3}\,(= 50^{\circ}\,57') = 230^{\circ}\,57'$\\
$r_{1} = 1{\cdot}0847$ & $r_{2} = 1{\cdot}0157$ & $r_{3} = 1{\cdot}0323.$\\
\end{tabular}
\end{center}

The calculation of the equation of the orbit is then as follows:
\[
\begin{aligned}
\lambda &= {\cdot}36397,\\
\lambda^{2} &= {\cdot}13248,\\
\lambda^{2} - 3 &= -2{\cdot}86752,\\
\sqrt{1 + \lambda^{2}} &= 1{\cdot}06477,\\
6\sqrt{1 + \lambda^{2}} &= 6{\cdot}38862,\\
2\sqrt{3}(\lambda^{2} - 3) &= -9{\cdot}93330,\\
2(\lambda^{2} - 3) &= -5{\cdot}73504,
\end{aligned}
\]
\[
R_{1} = 1{\cdot}0847,\quad R_{2} = +2{\cdot}4010,\quad R_{3} = -1{\cdot}6889,
\]
\[
4R_{1} = 4{\cdot}3388,\quad 2R_{1} = 2{\cdot}1694,
\]
\[
A = {\cdot}038982,\quad B = -{\cdot}014265,\quad C = +{\cdot}10464,
\]
\[
\varpi\,(= 20^{\circ}\,6') = 160^{\circ}\,52',\quad e = {\cdot}04151,\quad 1 - e^{2} = {\cdot}998277,\quad a = 1{\cdot}0481.
\]

\medskip
\noindent\textit{[p.~448]}

\medskip

\noindent\textit{[The page is occupied entirely by columnar arithmetic continuing the specimen calculation begun on p.~447. The principal intermediate values are reproduced above; the column totals yield the same $A = {\cdot}038982$, $B = -{\cdot}014265$, $C = {\cdot}10464$, and lead through $e^{2} = {\cdot}001723$, $1 - e^{2} = {\cdot}998277$, to $a = 1{\cdot}0481$.]}

\medskip
\noindent\textit{[p.~449]}

\medskip

\noindent\textit{[The greater part of this page is blank in the original.]}

\medskip

The calculation of the Times is similar to that for the first planogram, and requires no further illustration.

\medskip

The Table for Planogram No.~$2$ is as follows:

\medskip
\noindent\textit{[pp.~450--451]}

\medskip

\begin{center}
\textit{Planogram No.~$2,\ b = 180^{\circ},\ c = 0^{\circ}$ to $90^{\circ}$.}

\smallskip

\noindent\textit{[The original lays the table across two facing pages: a left half (p.~450) carrying columns $c$, $x_{1}$, $y_{1}$, $x_{2}$, $y_{2}$, $x_{3}$, $y_{3}$, $r_{1}$, $\phi_{1}$, $r_{2}$, $\phi_{2}$, $r_{3}$, $\phi_{3}$; and a right half (p.~451) carrying the Equation of Orbit ($r = Ax + By + C$), the eccentricity $e$, the angle $\varpi$, $a$, the chords $\gamma_{12}$, $\gamma_{23}$, $\gamma_{31}$, and the times $T_{12}$, $T_{23}$, $T_{31}$, repeating $c$ at the right margin. Heading note: ``The mark $\sim$ in any of the $T$ columns shows that the Time does not exist.'' The numerical entries are reproduced below as a single combined table; columns marked all $+$ or all $-$ retain the printed signs.]}
\end{center}

\smallskip

\noindent\resizebox{\textwidth}{!}{%
\begin{tabular}{l@{\hspace{0.3em}}r@{\hspace{0.3em}}r@{\hspace{0.3em}}r@{\hspace{0.3em}}r@{\hspace{0.3em}}r@{\hspace{0.3em}}r@{\hspace{0.3em}}r@{\hspace{0.3em}}r@{\hspace{0.3em}}r@{\hspace{0.3em}}r@{\hspace{0.3em}}r@{\hspace{0.3em}}r@{\hspace{0.3em}}r@{\hspace{0.3em}}r@{\hspace{0.3em}}r@{\hspace{0.3em}}r@{\hspace{0.3em}}r@{\hspace{0.3em}}r@{\hspace{0.3em}}r@{\hspace{0.3em}}r@{\hspace{0.3em}}r@{\hspace{0.3em}}r@{\hspace{0.3em}}r}
\hline
& $c$ & $x_{1}$ & $y_{1}$ & $x_{2}$ & $y_{2}$ & $x_{3}$ & $y_{3}$ & $r_{1}$ & $\phi_{1}$ & $r_{2}$ & $\phi_{2}$ & $r_{3}$ & $\phi_{3}$ & Eq.\ of Orbit & $e$ & $\varpi$ & $a$ & $\gamma_{12}$ & $\gamma_{23}$ & $\gamma_{31}$ & $T_{12}$ & $T_{23}$ & $T_{31}$\\
\hline
Circle & $0^{\circ}$ & $1{\cdot}000$ & ${\cdot}000$ & $-{\cdot}500$ & $+{\cdot}866$ & $-{\cdot}500$ & $-{\cdot}866$ & $1{\cdot}000$ & $0^{\circ}\,0'$ & $1{\cdot}000$ & $120^{\circ}\,0'$ & $1{\cdot}000$ & $240^{\circ}\,0'$ & ${\cdot}000{\cdot}000+1{\cdot}000$ & ${\cdot}000$ & ind.$^{\circ}$ & $1{\cdot}000$ & $1{\cdot}732$ & $1{\cdot}732$ & $1{\cdot}732$ & $1{\cdot}000$ & $1{\cdot}000$ & $1{\cdot}000$\\
 & $5$ & $1{\cdot}004$ & ${\cdot}051$ & ${\cdot}481$ & ${\cdot}878$ & $-{\cdot}525$ & ${\cdot}853$ & $1{\cdot}005$ & $2\,52$ & $1{\cdot}001$ & $118\,42$ & $1{\cdot}001$ & $238\,40$ & $+{\cdot}0101\,-{\cdot}0015\,1{\cdot}0104$ & ${\cdot}010$ & $171\,19$ & $1{\cdot}010$ & $1{\cdot}729$ & $1{\cdot}832$ & $1{\cdot}678$ & ${\cdot}987$ & $1{\cdot}106$ & ${\cdot}953$\\
 & $10$ & $1{\cdot}015$ & ${\cdot}102$ & ${\cdot}467$ & ${\cdot}889$ & ${\cdot}557$ & ${\cdot}838$ & $1{\cdot}020$ & $5\,44$ & $1{\cdot}004$ & $117\,41$ & $1{\cdot}006$ & $236\,40$ & ${\cdot}039\,{\cdot}014\,1{\cdot}046$ & ${\cdot}041$ & $160\,52$ & $1{\cdot}048$ & $1{\cdot}724$ & $1{\cdot}991$ & $1{\cdot}668$ & ${\cdot}956$ & $1{\cdot}316$ & ${\cdot}946$\\
 & $15$ & $1{\cdot}035$ & ${\cdot}156$ & ${\cdot}456$ & ${\cdot}900$ & ${\cdot}598$ & ${\cdot}821$ & $1{\cdot}047$ & $8\,30$ & $1{\cdot}009$ & $116\,54$ & $1{\cdot}016$ & $233\,30$ & & & & & & & & & & \\
 & $20$ & $1{\cdot}064$ & ${\cdot}210$ & ${\cdot}450$ & ${\cdot}910$ & ${\cdot}650$ & ${\cdot}802$ & $1{\cdot}085$ & $11\,10$ & $1{\cdot}016$ & $116\,19$ & $1{\cdot}032$ & $230\,30$ & ${\cdot}083\,{\cdot}061\,1{\cdot}126$ & ${\cdot}103$ & $143\,48$ & $1{\cdot}138$ & $1{\cdot}718$ & $2{\cdot}244$ & $1{\cdot}710$ & ${\cdot}924$ & $1{\cdot}777$ & ${\cdot}962$\\
 & $25$ & $1{\cdot}103$ & ${\cdot}269$ & ${\cdot}447$ & ${\cdot}921$ & ${\cdot}719$ & ${\cdot}778$ & $1{\cdot}136$ & $13\,43$ & $1{\cdot}024$ & $115\,55$ & $1{\cdot}060$ & $227\,30$ & & & & & & & & & & \\
Ellipses & $30$ & $1{\cdot}155$ & ${\cdot}333$ & ${\cdot}448$ & ${\cdot}931$ & ${\cdot}812$ & ${\cdot}749$ & $1{\cdot}202$ & $16\,6$ & $1{\cdot}033$ & $115\,42$ & $1{\cdot}104$ & $222\,30$ & ${\cdot}135\,{\cdot}209\,1{\cdot}317$ & ${\cdot}248$ & $122\,54$ & $1{\cdot}404$ & $1{\cdot}740$ & $2{\cdot}684$ & $1{\cdot}826$ & ${\cdot}878$ & $3{\cdot}238$ & ${\cdot}966$\\
 & $35$ & $1{\cdot}221$ & ${\cdot}404$ & ${\cdot}452$ & ${\cdot}941$ & ${\cdot}939$ & ${\cdot}710$ & $1{\cdot}286$ & $18\,19$ & $1{\cdot}044$ & $115\,40$ & $1{\cdot}178$ & $217\,30$ & & & & & & & & & & \\
 & $40$ & $1{\cdot}305$ & ${\cdot}484$ & ${\cdot}460$ & ${\cdot}951$ & $1{\cdot}125$ & ${\cdot}657$ & $1{\cdot}392$ & $20\,21$ & $1{\cdot}056$ & $116\,48$ & $1{\cdot}302$ & $210\,30$ & ${\cdot}161\,{\cdot}395\,1{\cdot}527$ & ${\cdot}426$ & $112\,12$ & $1{\cdot}867$ & $1{\cdot}805$ & $3{\cdot}055$ & $1{\cdot}925$ & & & \\
 & $45$ & $1{\cdot}414$ & ${\cdot}577$ & ${\cdot}471$ & ${\cdot}962$ & $1{\cdot}414$ & ${\cdot}577$ & $1{\cdot}527$ & $22\,12$ & $1{\cdot}071$ & $116\,6$ & $1{\cdot}528$ & $202\,30$ & ${\cdot}186\,{\cdot}815\,1{\cdot}972$ & ${\cdot}836$ & $102\,50$ & $6{\cdot}554$ & $2{\cdot}016$ & $3{\cdot}659$ & $2{\cdot}063$ & ${\cdot}878$ & $48{\cdot}60$ & ${\cdot}849$\\
Parab. & $50$ & $1{\cdot}556$ & ${\cdot}688$ & ${\cdot}487$ & ${\cdot}974$ & $1{\cdot}925$ & ${\cdot}440$ & $1{\cdot}701$ & $23\,51$ & $1{\cdot}088$ & $116\,35$ & $1{\cdot}975$ & $192\,30$ & ${\cdot}191\,{\cdot}982\,2{\cdot}140$ & $1{\cdot}000$ & $101\,14$ & $\infty$ & $2{\cdot}100$ & $3{\cdot}831$ & $2{\cdot}097$ & ${\cdot}879$ & $\infty$ & ${\cdot}820$\\
 & $51^{\circ}\,0'$ & $1{\cdot}589$ & ${\cdot}713$ & ${\cdot}491$ & ${\cdot}976$ & $2{\cdot}077$ & ${\cdot}400$ & $1{\cdot}741$ & $24\,9$ & $1{\cdot}093$ & $116\,43$ & $2{\cdot}115$ & $190\,30$ & ${\cdot}196\,1{\cdot}150\,2{\cdot}319$ & $1{\cdot}166$ & $99\,39$ & $6434$ & & & & & & \\
 & $52$ & $1{\cdot}624$ & ${\cdot}739$ & ${\cdot}495$ & ${\cdot}978$ & $2{\cdot}256$ & ${\cdot}353$ & $1{\cdot}787$ & $24\,28$ & $1{\cdot}097$ & $116\,51$ & $2{\cdot}283$ & $188\,30$ & & & & & & & & & & \\
 & $54$ & $1{\cdot}701$ & ${\cdot}795$ & ${\cdot}504$ & ${\cdot}984$ & $2{\cdot}729$ & ${\cdot}229$ & $1{\cdot}878$ & $25\,2$ & $1{\cdot}105$ & $117\,7$ & $2{\cdot}738$ & $184\,30$ & ${\cdot}203\,1{\cdot}719\,2{\cdot}898$ & $1{\cdot}720$ & $96\,44$ & $1{\cdot}481$ & & & & & & \\
Hyperbs. & $55$ & $1{\cdot}743$ & ${\cdot}824$ & ${\cdot}509$ & ${\cdot}986$ & $3{\cdot}049$ & ${\cdot}145$ & $1{\cdot}928$ & $25\,18$ & $1{\cdot}109$ & $117\,16$ & $3{\cdot}053$ & $182\,30$ & ${\cdot}207\,2{\cdot}187\,3{\cdot}366$ & $2{\cdot}192$ & $95\,25$ & ${\cdot}855$ & $2{\cdot}781$ & $4{\cdot}890$ & $2{\cdot}258$ & ${\cdot}895$ & $\sim$ & ${\cdot}665$\\
 & $56^{\circ}\,18'$ & $1{\cdot}802$ & ${\cdot}866$ & ${\cdot}515$ & ${\cdot}990$ & $3{\cdot}601$ & ${\cdot}000$ & $1{\cdot}999$ & $25\,39$ & $1{\cdot}116$ & $117\,30$ & $3{\cdot}601$ & $180\,30$ & ${\cdot}212\,3{\cdot}074\,4{\cdot}227$ & $3{\cdot}081$ & $93\,56$ & ${\cdot}498$ & & & & & & \\
 & $59$ & $1{\cdot}942$ & ${\cdot}961$ & ${\cdot}530$ & ${\cdot}997$ & $5{\cdot}786$ & $+{\cdot}566$ & $2{\cdot}166$ & $26\,20$ & $1{\cdot}129$ & $118\,59$ & $5{\cdot}813$ & $174\,30$ & ${\cdot}221\,-1{\cdot}415\,+15{\cdot}42$ & $1{\cdot}415$ & $90\,30$ & ${\cdot}077$ & & & & & & \\
Line & $60$ & $2{\cdot}000$ & $1{\cdot}000$ & ${\cdot}536$ & $1{\cdot}000$ & $7{\cdot}468$ & $1{\cdot}000$ & $2{\cdot}236$ & $26\,34$ & $1{\cdot}134$ & $118\,11$ & $7{\cdot}534$ & $172\,30$ & ${\cdot}224\,\infty\,\infty(y-1)$ & $\infty$ & $90$ & ${\cdot}000$ & $6{\cdot}932$ & $9{\cdot}468$ & $2{\cdot}536$ & ${\cdot}000$ & $\sim$ & ${\cdot}000$\\
 & $61$ & $2{\cdot}063$ & $1{\cdot}042$ & ${\cdot}542$ & $1{\cdot}003$ & $-10{\cdot}53$ & $+1{\cdot}793$ & $2{\cdot}311$ & $26\,48$ & $1{\cdot}140$ & $118\,24$ & $10{\cdot}68$ & $170\,30$ & & & & & Convex Orbits. & & & & & \\
Convex \{ & $63^{\circ}\,26'-\epsilon$ & $2{\cdot}236$ & $1{\cdot}155$ & ${\cdot}559$ & $1{\cdot}010$ & $\infty$ & $\infty$ & $2{\cdot}517$ & $27\,19$ & $1{\cdot}155$ & $118\,57$ & $\infty$ & $165\,30$ & ${\cdot}234\,+4{\cdot}906\,-3{\cdot}671$ & $4{\cdot}912$ & $87\,17$ & ${\cdot}159$ & & & & & & \\
 & $63^{\circ}\,26'+\epsilon$ & & & & & $-\infty$ & $\infty$ & & & & & $-\infty$ & $134\,30$ & ${\cdot}578\,-1{\cdot}761\,+3{\cdot}257$ & $1{\cdot}853$ & $108\,11$ & $1{\cdot}338$ & $\infty$ & $\infty$ & $2{\cdot}799$ & $\sim$ & $\infty$ & ${\cdot}909$\\
 & $64$ & $2{\cdot}281$ & $1{\cdot}184$ & ${\cdot}563$ & $1{\cdot}012$ & $+45{\cdot}22$ & $-12{\cdot}60$ & $2{\cdot}570$ & $27\,26$ & $1{\cdot}157$ & $119\,6$ & $46{\cdot}94$ & $344\,30$ & & & & & & & & & & \\
Hyperbs.\{ & $65$ & $2{\cdot}366$ & $1{\cdot}238$ & ${\cdot}571$ & $1{\cdot}015$ & $16{\cdot}36$ & $5{\cdot}146$ & $2{\cdot}670$ & $27\,37$ & $1{\cdot}165$ & $119\,21$ & $17{\cdot}15$ & $342\,30$ & ${\cdot}587\,{\cdot}979\,2{\cdot}134$ & $2{\cdot}494$ & $120\,21$ & $8{\cdot}666$ & $18{\cdot}014$ & $15{\cdot}380$ & $2{\cdot}945$ & & & \\
Parab. & $65^{\circ}\,52'$ & $2{\cdot}446$ & $1{\cdot}289$ & ${\cdot}578$ & $1{\cdot}019$ & $10{\cdot}12$ & $3{\cdot}552$ & $2{\cdot}765$ & $27\,47$ & $1{\cdot}171$ & $119\,35$ & $10{\cdot}80$ & $340\,30$ & ${\cdot}591\,{\cdot}805\,2{\cdot}257$ & $1{\cdot}000$ & $126\,19$ & $\infty$ & $11{\cdot}633$ & $9{\cdot}072$ & $3{\cdot}036$ & $\infty$ & $7{\cdot}746$ & $1{\cdot}386$\\
 & $66$ & $2{\cdot}459$ & $1{\cdot}297$ & ${\cdot}579$ & $1{\cdot}019$ & $9{\cdot}987$ & $3{\cdot}500$ & $2{\cdot}779$ & $27\,48$ & $1{\cdot}172$ & $119\,37$ & $10{\cdot}59$ & $340\,30$ & ${\cdot}593\,{\cdot}779\,2{\cdot}221$ & ${\cdot}979$ & $127\,15$ & $5{\cdot}383$ & & & & & & \\
 & $68$ & $2{\cdot}669$ & $1{\cdot}429$ & ${\cdot}596$ & $1{\cdot}026$ & $5{\cdot}617$ & $2{\cdot}369$ & $3{\cdot}028$ & $28\,10$ & $1{\cdot}186$ & $120\,11$ & $6{\cdot}090$ & $337\,30$ & ${\cdot}606\,{\cdot}338\,1{\cdot}894$ & ${\cdot}693$ & $150\,53$ & $3{\cdot}645$ & & & & & & \\
 & $70$ & $2{\cdot}924$ & $1{\cdot}586$ & ${\cdot}616$ & $1{\cdot}033$ & $3{\cdot}912$ & $1{\cdot}927$ & $3{\cdot}326$ & $28\,29$ & $1{\cdot}202$ & $120\,48$ & $4{\cdot}360$ & $335\,30$ & ${\cdot}619\,-{\cdot}120\,1{\cdot}708$ & ${\cdot}630$ & $169$ & $2{\cdot}834$ & $5{\cdot}409$ & $3{\cdot}649$ & $3{\cdot}584$ & $5{\cdot}735$ & $6{\cdot}343$ & $2{\cdot}685$\\
Ellipses & $72$ & $3{\cdot}237$ & $1{\cdot}777$ & ${\cdot}638$ & $1{\cdot}041$ & $3{\cdot}008$ & $1{\cdot}694$ & $3{\cdot}693$ & $28\,47$ & $1{\cdot}221$ & $121\,29$ & $3{\cdot}455$ & $330\,30$ & ${\cdot}635\,+{\cdot}027\,1{\cdot}599$ & ${\cdot}636$ & $182\,25$ & $2{\cdot}674$ & & & & & & \\
 & $75$ & $3{\cdot}864$ & $2{\cdot}155$ & ${\cdot}674$ & $1{\cdot}054$ & $2{\cdot}230$ & $1{\cdot}488$ & $4{\cdot}424$ & $29\,9$ & $1{\cdot}251$ & $122\,36$ & $2{\cdot}681$ & $325\,30$ & ${\cdot}654\,{\cdot}185\,1{\cdot}497$ & ${\cdot}680$ & $195\,47$ & $2{\cdot}783$ & $3{\cdot}859$ & $3{\cdot}981$ & $4{\cdot}644$ & $2{\cdot}62$ & $6{\cdot}68$ & $4{\cdot}64$\\
 & $80$ & $5{\cdot}759$ & $3{\cdot}274$ & ${\cdot}751$ & $1{\cdot}079$ & $1{\cdot}568$ & $1{\cdot}312$ & $6{\cdot}624$ & $29\,37$ & $1{\cdot}315$ & $124\,49$ & $2{\cdot}045$ & $320\,30$ & ${\cdot}692\,{\cdot}366\,1{\cdot}439$ & ${\cdot}783$ & $207\,52$ & $3{\cdot}716$ & $3{\cdot}327$ & $6{\cdot}212$ & $6{\cdot}870$ & $1{\cdot}97$ & $10{\cdot}21$ & $9{\cdot}343$\\
 & $85$ & $11{\cdot}47$ & $6{\cdot}599$ & ${\cdot}854$ & $1{\cdot}112$ & $1{\cdot}217$ & $1{\cdot}216$ & $13{\cdot}25$ & $29\,54$ & $1{\cdot}402$ & $127\,32$ & $1{\cdot}720$ & $315\,30$ & ${\cdot}740\,{\cdot}514\,1{\cdot}455$ & ${\cdot}892$ & $214\,47$ & $7{\cdot}721$ & $3{\cdot}115$ & $12{\cdot}895$ & $13{\cdot}480$ & & & \\
Parab. & $89^{\circ}\,20'$ & $86{\cdot}41$ & $49{\cdot}79$ & ${\cdot}979$ & $1{\cdot}148$ & $1{\cdot}024$ & $1{\cdot}161$ & $99{\cdot}5$ & $29\,56$ & $1{\cdot}508$ & $130\,24$ & $1{\cdot}548$ & $311\,30$ & ${\cdot}764\,{\cdot}645\,1{\cdot}505$ & $1{\cdot}000$ & $219\,51$ & $\infty$ & $3{\cdot}055$ & $9{\cdot}943$ & $102{\cdot}7$ & $1{\cdot}200$ & $225{\cdot}4$ & $\infty$\\
Hyperbs. & $90-\epsilon$ & $\infty$ & $\infty$ & $1{\cdot}000$ & $1{\cdot}155$ & $+1{\cdot}000$ & $-1{\cdot}155$ & $\infty$ & $30$ & $1{\cdot}527$ & $130\,54$ & $1{\cdot}527$ & $310\,30$ & ${\cdot}770\,+{\cdot}666\,+1{\cdot}527$ & $1{\cdot}018$ & $220\,6$ & $41{\cdot}000$ & $3{\cdot}055$ & $\infty$ & $\infty$ & $1{\cdot}148$ & $\infty$ & $\sim$\\
\hline
\end{tabular}}

\medskip

\noindent\textit{[p.~452]}

\medskip

\begin{center}
\textit{Article Nos.~$95$ to $98$. Planogram No.~$3$, the Orbit-pole at one of the points $A$.}
\end{center}

\paragraph{95.} When the orbit-pole is at one of the points $A$, the orbit-plane passes through one of the rays, and as there is no longer on this ray any determinate point of intersection, the orbit (as was seen) becomes indeterminate. Thus consider the point $A$ for which $b = 270^{\circ}$, $c = 60^{\circ}$: we have
\[
\begin{aligned}
\alpha,\ \beta,\ \gamma &= -1,\ 0,\ 0,\\
\alpha',\ \beta',\ \gamma' &= 0,\ -\tfrac{1}{2},\ -\tfrac{1}{2}\sqrt{3},\\
\alpha'',\, \beta'',\, \gamma'' &= 0,\ -\tfrac{1}{2}\sqrt{3},\ \tfrac{1}{2},
\end{aligned}
\]

\begin{center}
\textit{[Figure 9: a circle in the $(x,y)$-plane of unit radius, with the four cardinal positions of points $1$, $2$, $3$, $B$ marked. Point $1$ at the top, $2$ at lower left, $3$ at lower right, $B$ at the right.]}
\end{center}

and consequently the formula gives
\[
\begin{aligned}
x_{1}'\,:\,y_{1}'\,:\,1 &= 0\ :\ 0\ :\ 0,\\
x_{2}'\,:\,y_{2}'\,:\,1 &= -\tfrac{1}{2}\sqrt{3} - \sqrt{3}\ :\ 3\ :\ -\tfrac{1}{2}\sqrt{3} - \sqrt{3},\\
x_{3}'\,:\,y_{3}'\,:\,1 &= -\tfrac{1}{2}\sqrt{3} - \sqrt{3}\ :\ -3\ :\ -\tfrac{1}{2}\sqrt{3} - \sqrt{3},
\end{aligned}
\]
and, moreover, $x = -x'$, $y = -y'$. From the formula the value of $x_{1}'$ or $x_{1}$ is given as $\tfrac{0}{0}$, but the true value is obviously $x_{1} = 1$; the value of $y_{1}$ is actually indeterminate. The formulae give the values of $(x_{2},\, y_{2})$, $(x_{3},\, y_{3})$, viz.\ the system is
\[
\begin{aligned}
x_{1} &= 1, & y_{1} &= \text{ind.}\\
x_{2} &= -1, & y_{2} &= \tfrac{2}{\sqrt{3}}, & \text{whence } r_{2} = r_{3} = \sqrt{\tfrac{7}{3}},\\
x_{3} &= -1, & y_{3} &= -\tfrac{2}{\sqrt{3}},
\end{aligned}
\]

\medskip
\noindent\textit{[p.~453]}

\medskip

so that the orbits in the planogram are the whole series of conics having a given focus, $S$, and passing through two fixed points, $2,\, 3$, having the common abscissa $x = -1$, and at equal distances $\tfrac{2}{\sqrt{3}}$ ($= 1{\cdot}15470$) on opposite sides of the axis. The axis of $x$ is obviously the common transverse axis for all the orbits; that is, the equation of the orbit will be of the form $r = Ax + B$; and writing $x = -1$, we have $\sqrt{\tfrac{7}{3}} = -A + B$, viz.\ the equation is $r - \sqrt{\tfrac{7}{3}} = A(x + 1)$; the value of $A$ will be determined if we assume for the point $1$ a determinate position on the line $x = 1$, say its ordinate is $= y_{1}$; for then if $r_{1} = \sqrt{1 + y_{1}^{2}}$ we have $r_{1} - \sqrt{\tfrac{7}{3}} = 2A$, and the equation is $r - \sqrt{\tfrac{7}{3}} = \tfrac{1}{2}(r_{1} - \sqrt{\tfrac{7}{3}})(x + 1)$. In particular if $y_{1} = 0$, we have $r_{1} = 1$, and the equation of the orbit is $r - \sqrt{\tfrac{7}{3}} = \tfrac{1}{2}(1 - \sqrt{\tfrac{7}{3}})(x + 1)$: this is the orbit, eccentricity $= \tfrac{1}{2}(\sqrt{\tfrac{7}{3}} - 1) = {\cdot}264$, belonging to the point $A$ as a point in planogram No.~$1$: for the value of $y_{1}$ being in that planogram originally assumed $= 0$, is of course $= 0$ when the orbit-pole comes to be the point $A$.

\paragraph{96.} We may conversely take the equation of the orbit, or say the value of $A\,(= \pm e)$ in the equation $r - \sqrt{\tfrac{7}{3}} = A(x + 1)$, to be given; and then writing $x = x_{1} = 1$, we have
\[
r_{1} = \sqrt{\tfrac{7}{3}} + 2A,\quad\text{that is}\quad y_{1}^{2} = (\sqrt{\tfrac{7}{3}} + 2A)^{2} - 1;
\]
for
\[
r_{1} = 1\quad\text{or}\quad y_{1} = 0,\quad A = \tfrac{1}{2}(1 - \sqrt{\tfrac{7}{3}}) = -{\cdot}264,
\]
and as $r_{1}$ increases to $r_{1} = \sqrt{\tfrac{7}{3}}$, or $y_{1}$ increases to $+\tfrac{2}{\sqrt{3}}$, $A$ diminishes from $-{\cdot}264$ to $0$; viz., for $r_{1} = \sqrt{\tfrac{7}{3}}$ or $y_{1} = \tfrac{2}{\sqrt{3}}$, the orbit is a circle; as $r_{1}$ increases from $\sqrt{\tfrac{7}{3}}$, or $y_{1}$ from $+\tfrac{2}{\sqrt{3}}$, $A$ increases from $0$ positively; for $r_{1} = \sqrt{\tfrac{7}{3}} + 2 = 3{\cdot}527$, or $y_{1} = +\tfrac{1}{\sqrt{3}}\sqrt{} = 2{\cdot}896$, $A$ becomes $= 1$; that is, the orbit is a parabola; and for larger positive values of $r_{1}$, or positive or negative values of $y_{1}$, the orbit is a hyperbola (concave); and ultimately for $r_{1} = \infty$ or $y_{1} = \pm\infty$, the orbit is the right line $x + 1 = 0$. Thus $A$ extends from $-{\cdot}264$ to $0$, and thence from $0$ positively to $\infty$.

\paragraph{97.} In further illustration, suppose that the orbit-pole, instead of being at $A$, is a point in the immediate neighbourhood of $A$, say that the rectangular spherical coordinates, measured from $A$ in the direction of the meridian and perpendicular thereto, are $\xi$ and $\eta$, the colatitude and longitude of the orbit-pole being thus $c = 60^{\circ} + \xi$, and $b = 270^{\circ} + \dfrac{2}{\sqrt{3}}\eta$; we have then, $\xi,\, \eta$ being indefinitely small,
\[
\begin{aligned}
\alpha,\ \beta,\ \gamma &= -1,\ -\tfrac{2}{\sqrt{3}}\eta,\ 0,\\
\alpha',\ \beta',\ \gamma' &= \tfrac{2}{\sqrt{3}}\eta\,,\ -\tfrac{1}{2} - \tfrac{\sqrt{3}}{2}\xi\,,\ -\tfrac{\sqrt{3}}{2} + \tfrac{1}{2}\xi\,,\\
\alpha'',\, \beta'',\, \gamma'' &= \eta\,,\ -\tfrac{\sqrt{3}}{2} + \tfrac{1}{2}\xi\,,\ \tfrac{1}{2} - \tfrac{\sqrt{3}}{2}\xi\,.
\end{aligned}
\]

\end{document}
